% comp 06: summaries; phase 2 rewrite (the opening justification sentence of the old SPLIT paragraph is dropped; comp/s04 now says what the proposition gives)
\section{What the network sees: summaries}\label{comp:sec:summaries}

The network never sees the data. It sees $S(D)$, a vector we chose, and learns $\pi(\theta\mid S=s)$ (Proposition~\ref{gbc:prop:pinball}). A summary\index{summary statistics|textbf} is a promise about what we are willing to forget. If $S$ is sufficient the promise costs nothing; if not, this valid posterior conditions on less than the data and is less concentrated on average over data sets. Summarizing is the first of two approximations to $\pi(\theta\mid D^{\mathrm{obs}})$; the second is finite training.

Sufficient summaries lose nothing, as in the normal mean and the two arm trial of Section~\ref{gbc:sec:qnnexplained}, but when the support does not depend on $\theta$, a sufficient statistic of fixed dimension exists, under regularity conditions, only for exponential families. Learned summaries start from Fearnhead and Prangle (2012), who take $S(D)=\E[\theta\mid D]$ and estimate it by regressing simulated $\theta$ on simulated data. By the tower property the posterior given this summary keeps the full posterior mean, and its variance is the average of the full posterior variances over data sets sharing the summary. 

The posterior mean can carry no information. If the prior and the likelihood are unchanged when $\theta$ is replaced by $-\theta$, the posterior is symmetric about zero, typically with a mode on each side, and $\E[\theta\mid D]=0$ for every $D$, so the Fearnhead and Prangle summary is a constant and the network returns the prior. Coverage on held out simulations does not reveal the failure, since intervals from $\pi(\theta\mid S)$ have their nominal coverage over the prior predictive distribution\index{prior predictive distribution|textbf} for every $S$, and the global and local ranks of Section~\ref{comp:sec:checks} are flat for a network that returns the prior. The held out loss against a network given another summary shows it, and so does a reference posterior at the observed summary.

The third kind puts in what a frequentist would compute. For a binary treatment a natural choice collects the weighted outcome means, the covariate sums by arm divided by $N$ (the arm means times the arm shares) and the covariate second moments,
\[
\begin{split}
S(D) = \Bigl(&\tfrac1N\textstyle\sum_i \tfrac{W_iY_i}{\hat e(X_i)},\ \tfrac1N\sum_i\tfrac{(1-W_i)Y_i}{1-\hat e(X_i)},\\
&\tfrac1N\textstyle\sum_i X_iW_i,\ \tfrac1N\sum_i X_i(1-W_i),\ \tfrac1N\sum_i X_iX_i^\top\Bigr),
\end{split}
\]
to which one adds outcome regression coefficients by arm, treated and control means within propensity strata and any consistent estimator of the target. The first two components are inverse probability weighted outcome means (Chapter~\ref{ch:ipw}), whose difference is the Horvitz and Thompson estimate of the average effect. The covariate sums record imbalance; the second moments let the network learn its importance for the outcome model.

Fit $\hat e(X)$ inside $S$ on every simulated data set. If the observed analysis cross fits its nuisances, repeat that cross fit inside every simulated $S$, at the cost of nuisance fits on all $J$ data sets, so the network learns the law of the statistic it will see.

Chapter~\ref{ch:blp} proves what a good summary buys (Proposition~\ref{blp:prop:bvm}). Under regularity conditions the posterior given $S$ is asymptotically normal, centred at the efficient minimum distance estimator built from $S$. If $S$ contains an asymptotically normal estimator of a subvector of $\theta$, the summary posterior for that subvector is asymptotically at least as concentrated as the estimator, with a limiting covariance no larger in the matrix order. A good estimator of the target inside $S$ therefore bounds the loss from summarizing. The conditions include a smooth injective limit of $S$, a prior density continuous and positive at the truth, and conditions on the density of $S$, because convergence in distribution of $S$ does not by itself give convergence of the posterior given $S$. The result concerns the exact summary posterior, not a trained network.

No summary of a covariate can produce heterogeneity in $\tau(x)$ if the outcome model omits its interaction with treatment. Choose the summary with the model.

The second approximation is finite training: the network class, the optimizer and the number of simulations $J$. Padilla, Tansey and Chen (2022) give rates for quantile regression with ReLU networks, for an exact empirical risk minimizer at one fixed quantile level; for targets with compositional structure, an intrinsic dimension replaces the summary dimension $k$. Our network is trained at levels drawn at random, for which we know no rate. Theorem A.4 of Polson and Sokolov (2026a), proved there in outline, gives consistency for fully connected ReLU networks of the level and the summary: for a bounded parameter, summaries in a compact set, a H\"older continuous conditional quantile function and a width that grows slowly enough with $J$, the expected integrated check loss of the empirical minimizer converges to its infimum, the expectation being over the training set as well. It does not cover the product architecture of Section~\ref{comp:sec:network}, whose complexity we have not bounded, and none of it bounds the error of the stochastic gradient descent that trains the network.
